\documentclass[11pt]{article}
\usepackage[T1]{fontenc}
\usepackage[spanish]{babel}
\usepackage{amsmath,amssymb}
\usepackage[a4paper,margin=2.3cm]{geometry}
\usepackage{graphicx}
\usepackage{tikz}
\usetikzlibrary{arrows.meta,patterns.meta}

\IfFileExists{cilindro_hueco_visualizacion.pdf}{%
  \newcommand{\includevisualizacion}[2][]{\includegraphics[##1]{##2.pdf}}%
}{%
  \usepackage[inkscapelatex=false]{svg}
  \newcommand{\includevisualizacion}[2][]{\includesvg[##1]{##2}}%
}

\spanishdecimal{.}
\setlength{\parindent}{0pt}
\setlength{\parskip}{0.35em}
\setlength{\abovedisplayskip}{4pt}
\setlength{\belowdisplayskip}{4pt}
\setlength{\abovedisplayshortskip}{4pt}
\setlength{\belowdisplayshortskip}{4pt}

\newcommand{\encabezado}[3]{%
  {\Large\bfseries #1\par}
  \vspace{2pt}
  {\normalsize #2\hfill #3\par}
  \vspace{6pt}\hrule\vspace{8pt}}

\newcommand{\problema}[2]{%
  \vspace{0.55em}
  \noindent\textbf{#1.}\enspace{\small #2}\par
  \vspace{0.3em}}

\newcommand{\inciso}[1]{%
  \vspace{0.35em}
  \noindent\textbf{(#1)}\enspace}

\begin{document}

\encabezado{Campo y potencial electrostático de un cilindro hueco infinito}
{Electrodinámica I}{P.S., J.A., J.R.}

\problema{1}{Considere un cilindro muy largo (aproximado como infinito) de radio interior $a$ y radio exterior $b$. El cilindro posee una densidad de carga volumétrica no uniforme que depende de la distancia radial $r$ al eje del cilindro, descrita por:
\begin{equation}
\rho(r) =
\begin{cases}
0, & 0\leq r<a,\\[2pt]
\dfrac{\rho_0}{r}, & a<r<b,\\[5pt]
0, & r>b,
\end{cases}
\end{equation}
donde $\rho_0$ es una constante con unidades de carga por unidad de área. 
Se pide:
\begin{enumerate}
    \item Calcular el campo eléctrico $\vec E(r)$ en todo el espacio utilizando integración directa (ley de Coulomb).
    \item Obtener el mismo campo eléctrico empleando la Ley de Gauss en forma integral.
    \item Calcular el potencial eléctrico $\phi(r)$ en todo el espacio, estableciendo como condición de referencia $\phi(b) = 0$.
\end{enumerate}
}

\begin{figure}[htbp]
\centering
\begin{tikzpicture}[
    scale=0.74,
    transform shape,
    line cap=round,
    line join=round,
    eje/.style={-{Stealth[length=1.8mm]},semithick},
    cota/.style={{Stealth[length=1.3mm]}-{Stealth[length=1.3mm]}}
]
\begin{scope}[x={(1cm,0cm)},y={(0.43cm,0.21cm)},z={(0cm,1cm)}]
\def\raen{0.92}
\def\rben{1.82}
\def\hen{3.35}

% Región cargada visible en la cara superior y en dos generatrices.
\path[
  pattern={Lines[angle=45,distance=3.1pt,line width=0.24pt]},
  pattern color=black!60,
  even odd rule
]
  plot[domain=0:360,samples=90,variable=\t]
       ({\rben*cos(\t)},{\rben*sin(\t)},\hen)--cycle
  plot[domain=0:360,samples=90,variable=\t]
       ({\raen*cos(\t)},{\raen*sin(\t)},\hen)--cycle;
\path[
  pattern={Lines[angle=45,distance=3.1pt,line width=0.24pt]},
  pattern color=black!60
] (\raen,0,0)--(\rben,0,0)--(\rben,0,\hen)--(\raen,0,\hen)--cycle;
\path[
  pattern={Lines[angle=45,distance=3.1pt,line width=0.24pt]},
  pattern color=black!60
] (-\raen,0,0)--(-\rben,0,0)--(-\rben,0,\hen)--(-\raen,0,\hen)--cycle;

% Contornos del cascarón cilíndrico.
\draw[thick] plot[domain=0:360,samples=90,variable=\t]
  ({\rben*cos(\t)},{\rben*sin(\t)},\hen);
\draw[thick] plot[domain=0:360,samples=90,variable=\t]
  ({\raen*cos(\t)},{\raen*sin(\t)},\hen);
\draw[thick] plot[domain=180:360,samples=45,variable=\t]
  ({\rben*cos(\t)},{\rben*sin(\t)},0);
\draw[densely dashed,black!55] plot[domain=0:180,samples=45,variable=\t]
  ({\rben*cos(\t)},{\rben*sin(\t)},0);
\draw[thick] plot[domain=180:360,samples=45,variable=\t]
  ({\raen*cos(\t)},{\raen*sin(\t)},0);
\draw[densely dashed,black!55] plot[domain=0:180,samples=45,variable=\t]
  ({\raen*cos(\t)},{\raen*sin(\t)},0);
\foreach \r in {-\rben,-\raen,\raen,\rben}{
  \draw[thick] (\r,0,0)--(\r,0,\hen);
  \draw[densely dashed,black!55] (\r,0,-0.60)--(\r,0,0);
  \draw[densely dashed,black!55] (\r,0,\hen)--(\r,0,{\hen+0.60});
}

% Eje y radios característicos.
\coordinate (Oen) at (0,0,0);
\draw[eje] (0,0,-0.72)--(0,0,{\hen+0.92}) node[left] {$z$};
\fill (Oen) circle (1.1pt);
\draw[cota] (Oen)--({\raen*cos(218)},{\raen*sin(218)},0)
  node[pos=.64,above left=-1pt] {$a$};
\draw[cota] (Oen)--({\rben*cos(252)},{\rben*sin(252)},0)
  node[pos=.72,below left=-1pt] {$b$};
\node[anchor=west,align=left] at (2.32,0,2.25)
  {$\displaystyle \rho(r)=\frac{\rho_0}{r}$\\[-1pt]$a<r<b$};
\draw[-{Stealth[length=1.5mm]}] (2.22,0,2.06)--(1.45,0,1.83);
\node[anchor=west,font=\scriptsize] at (2.30,0,3.35)
  {invariante bajo traslaciones en $z$};
\end{scope}
\end{tikzpicture}

\caption{Cascarón cilíndrico infinito. El tramado identifica la región cargada
$a<r<b$, donde $\rho(r)=\rho_0/r$.}
\label{fig:geometria-cilindro}
\end{figure}

\textbf{Solución:} La distribución es invariante bajo rotaciones alrededor del eje
$z$ y bajo traslaciones a lo largo de ese eje. Por esto el campo sólo puede depender
de $r$ y apuntar en dirección radial:
\begin{equation*}
\vec E(r)=E(r)\,\hat{r}.
\end{equation*}

\inciso{a} Para calcular el campo mediante integración directa partimos de la ley
de Coulomb para una distribución continua. Conservamos la convención según la cual $\vec{x}$ representa el punto de observación y
$\vec{x}'$ un punto fuente, escribiendo la ley directamente en
forma vectorial:
\begin{equation}
\vec E(\vec{x})=
\frac{1}{4\pi\varepsilon_0}
\int_V
\rho(\vec{x}')
\frac{\vec{x}-\vec{x}'}{\lvert\vec{x}-\vec{x}'\rvert^3}\,dV'.
\end{equation}

\paragraph{Elección de coordenadas y variables.}
Usamos un único sistema de coordenadas cilíndricas, cuyo eje es el eje $z$ del
cilindro. El punto de observación $P$ permanece fijo durante la integración y, por
simetría, podemos escogerlo sobre el semieje $x$ positivo. El punto fuente $Q$, en
cambio, recorre el volumen cargado:
\[
\begin{aligned}
P:&\quad \vec{x}=r\,\hat{x}
&&\Longleftrightarrow &&(r,\varphi,z)=(r,0,0),\\
Q:&\quad \vec{x}'=\vec{x}'(r',\varphi',z')
&&\Longleftrightarrow &&(r',\varphi',z').
\end{aligned}
\]
Así, $r$ es el radio donde se desea conocer el campo y actúa como parámetro: no se
integra respecto de él. Por el contrario, $r'$, $\varphi'$ y $z'$ son variables de
integración, con $a<r'<b$, $0\leq\varphi'<2\pi$ y
$-\infty<z'<\infty$. Cada posición de $Q$ determina el valor local de la densidad
de carga,
\[
\rho(\vec{x}')=\rho(r')=\frac{\rho_0}{r'}.
\]

El elemento de volumen de la fuente en coordenadas cilíndricas es
\begin{equation}
dV'=r'\,dr'\,d\varphi'\,dz'.
\end{equation}
La posición del elemento de carga es
$\vec{x}'=r'\cos\varphi'\,\hat{x}
+r'\sin\varphi'\,\hat{y}+z'\,\hat{z}$, y el vector de separación es
\begin{equation}
\vec{x}-\vec{x}'=
(r-r'\cos\varphi')\hat{x}
-r'\sin\varphi'\hat{y}-z'\hat{z}.
\end{equation}

La figura~\ref{fig:integracion-directa} muestra un segmento del cilindro infinito,
el punto fuente, el punto de observación y los vectores que intervienen en la ley de
Coulomb. El tramado se usa únicamente para distinguir la región que contiene carga.

\begin{figure}[htbp]
\centering
\begin{tikzpicture}[
    scale=0.90,
    transform shape,
    line cap=round,
    line join=round,
    eje/.style={-{Stealth[length=2.0mm]},semithick},
    vector/.style={-{Stealth[length=2.4mm]},very thick},
    auxiliar/.style={densely dashed,black!55},
    cota/.style={{Stealth[length=1.5mm]}-{Stealth[length=1.5mm]}},
    every node/.style={font=\small}
]
% La proyección oblicua se usa sólo para representar la geometría 3D.
\begin{scope}[
    x={(1cm,0cm)},
    y={(0.44cm,0.22cm)},
    z={(0cm,1cm)}
]
\def\ra{1.08}
\def\rb{2.22}
\def\altura{5.15}
\def\rp{1.68}
\def\varphip{38}
\def\zp{3.05}

\coordinate (O) at (0,0,0);
\coordinate (Zp) at (0,0,\zp);
\coordinate (Q) at ({\rp*cos(\varphip)},{\rp*sin(\varphip)},\zp);
\coordinate (P) at (3.35,0,0);

% Regiones cargadas visibles en un corte meridiano y en la cara superior.
% El tramado, y no un color sólido, identifica a<r'<b.
\path[
  pattern={Lines[angle=45,distance=3.2pt,line width=0.25pt]},
  pattern color=black!60
] (\ra,0,0)--(\rb,0,0)--(\rb,0,\altura)--(\ra,0,\altura)--cycle;
\path[
  pattern={Lines[angle=45,distance=3.2pt,line width=0.25pt]},
  pattern color=black!60
] (-\ra,0,0)--(-\rb,0,0)--(-\rb,0,\altura)--(-\ra,0,\altura)--cycle;
\path[
  pattern={Lines[angle=45,distance=3.2pt,line width=0.25pt]},
  pattern color=black!60,
  even odd rule
]
  plot[domain=0:360,samples=100,variable=\t]
       ({\rb*cos(\t)},{\rb*sin(\t)},\altura)--cycle
  plot[domain=0:360,samples=100,variable=\t]
       ({\ra*cos(\t)},{\ra*sin(\t)},\altura)--cycle;

% Contornos del cilindro hueco.
\draw[thick]
  plot[domain=0:360,samples=100,variable=\t]
       ({\rb*cos(\t)},{\rb*sin(\t)},\altura);
\draw[thick]
  plot[domain=0:360,samples=100,variable=\t]
       ({\ra*cos(\t)},{\ra*sin(\t)},\altura);
\draw[thick]
  plot[domain=180:360,samples=55,variable=\t]
       ({\rb*cos(\t)},{\rb*sin(\t)},0);
\draw[auxiliar]
  plot[domain=0:180,samples=55,variable=\t]
       ({\rb*cos(\t)},{\rb*sin(\t)},0);
\draw[thick]
  plot[domain=180:360,samples=55,variable=\t]
       ({\ra*cos(\t)},{\ra*sin(\t)},0);
\draw[auxiliar]
  plot[domain=0:180,samples=55,variable=\t]
       ({\ra*cos(\t)},{\ra*sin(\t)},0);

\foreach \x in {-\rb,\rb}{
  \draw[thick] (\x,0,0)--(\x,0,\altura);
}
\foreach \x in {-\ra,\ra}{
  \draw[black!65] (\x,0,0)--(\x,0,\altura);
}

% Sistema de referencia.
\draw[eje] (O)--(3.85,0,0) node[below] {$x$};
\draw[eje] (O)--(0,3.35,0) node[above right=-1pt] {$y$};
\draw[eje] (O)--(0,0,5.95) node[left] {$z$};
\fill (O) circle (1.1pt);
\node[below left] at (O) {$O$};

% Radios a y b en la base.
\draw[cota]
  (O)--({\ra*cos(218)},{\ra*sin(218)},0)
  node[pos=.62,above left=-1pt] {$a$};
\draw[cota]
  (O)--({\rb*cos(250)},{\rb*sin(250)},0)
  node[pos=.72,below left=-1pt] {$b$};

% Vector posición del punto de observación.
\draw[vector] (O)--(P)
  node[pos=.72,below=2pt,fill=white,inner sep=1pt]
  {$\vec{x}=r\,\hat{x}$};
\draw[thick,fill=white] (P) circle (2pt);
\node[above right=-1pt] at (P) {$P$};

% Coordenadas del punto fuente.
\draw[auxiliar] (Zp)--(2.30,0,\zp);
\draw[vector] (Zp)--(Q)
  node[pos=.55,above left=-1pt,fill=white,inner sep=1pt] {$r'$};
\coordinate (PhiStart) at (0.95,0,\zp);
\coordinate (PhiCOne) at ({0.98*cos(12)},{0.98*sin(12)},\zp);
\coordinate (PhiCTwo) at ({0.98*cos(27)},{0.98*sin(27)},\zp);
\coordinate (PhiEnd) at ({0.95*cos(\varphip)},{0.95*sin(\varphip)},\zp);
\draw[-{Stealth[length=1.5mm]}]
  (PhiStart) .. controls (PhiCOne) and (PhiCTwo) .. (PhiEnd);
\node[below=4pt] at ({0.98*cos(20)},{0.98*sin(20)},\zp) {$\varphi'$};
\draw[cota] (-0.26,0,0)--(-0.26,0,\zp)
  node[midway,left] {$z'$};

% Vectores x' y x-x'.
\draw[vector,densely dashed] (O)--(Q)
  node[pos=.58,left=2pt,fill=white,inner sep=1pt] {$\vec{x}'$};
\draw[vector,line width=1.15pt] (Q)--(P)
  node[pos=.55,right=2pt,fill=white,inner sep=1pt]
  {$\vec{x}-\vec{x}'$};
\fill (Q) circle (2.2pt);
\node[above right=1pt] at (Q) {$Q$};

% Densidad y significado del tramado.
\node[anchor=east,align=left,fill=white,inner sep=2pt]
  at (-2.75,0,4.15)
  {$\displaystyle\rho(r')=\frac{\rho_0}{r'}$\\[-1pt]$a<r'<b$};
\draw[-{Stealth[length=1.5mm]}]
  (-2.65,0,3.85)--(-1.72,0,3.55);
\end{scope}

% Ampliación del elemento cilíndrico: una cuña anular, no una caja cartesiana.
\begin{scope}[shift={(7.15,3.20)}]
  \def\ri{0.76}      % radio interior de la cuña ampliada
  \def\ro{1.36}      % radio exterior de la cuña ampliada
  \def\anga{-32}     % ángulo inicial
  \def\angb{42}      % ángulo final
  \def\cx{-0.52}     % centro de curvatura local
  \def\cy{-0.03}
  \def\ex{0.28}      % desplazamiento que representa dz'
  \def\ey{0.48}

  \coordinate (Ov)  at (\cx,\cy);
  \coordinate (Aib) at ({\cx+\ri*cos(\anga)},{\cy+\ri*sin(\anga)});
  \coordinate (Aob) at ({\cx+\ro*cos(\anga)},{\cy+\ro*sin(\anga)});
  \coordinate (Bob) at ({\cx+\ro*cos(\angb)},{\cy+\ro*sin(\angb)});
  \coordinate (Bib) at ({\cx+\ri*cos(\angb)},{\cy+\ri*sin(\angb)});
  \coordinate (Ait) at ({\cx+\ex+\ri*cos(\anga)},{\cy+\ey+\ri*sin(\anga)});
  \coordinate (Aot) at ({\cx+\ex+\ro*cos(\anga)},{\cy+\ey+\ro*sin(\anga)});
  \coordinate (Bot) at ({\cx+\ex+\ro*cos(\angb)},{\cy+\ey+\ro*sin(\angb)});
  \coordinate (Bit) at ({\cx+\ex+\ri*cos(\angb)},{\cy+\ey+\ri*sin(\angb)});

  % Cara curva exterior y cara radial visible.
  \path[
    pattern={Dots[distance=3pt,radius=0.24pt]},
    pattern color=black!48
  ]
    (Aob) arc[start angle=\anga,end angle=\angb,radius=\ro]
    --(Bot) arc[start angle=\angb,end angle=\anga,radius=\ro]--cycle;
  \draw[thick]
    (Aob) arc[start angle=\anga,end angle=\angb,radius=\ro]
    --(Bot) arc[start angle=\angb,end angle=\anga,radius=\ro]--cycle;

  \path[
    pattern={Dots[distance=3pt,radius=0.24pt]},
    pattern color=black!48
  ] (Bob)--(Bib)--(Bit)--(Bot)--cycle;
  \draw[thick] (Bob)--(Bib)--(Bit)--(Bot)--cycle;

  % Cara superior: sector anular limitado por dos arcos y dos lados radiales.
  \path[
    pattern={Lines[angle=45,distance=3.2pt,line width=0.22pt]},
    pattern color=black!58
  ]
    (Ait)--(Aot)
    arc[start angle=\anga,end angle=\angb,radius=\ro]
    --(Bit)
    arc[start angle=\angb,end angle=\anga,radius=\ri]--cycle;
  \draw[very thick]
    (Ait)--(Aot)
    arc[start angle=\anga,end angle=\angb,radius=\ro]
    --(Bit)
    arc[start angle=\angb,end angle=\anga,radius=\ri]--cycle;

  % Aristas restantes de la extrusión dz'.
  \draw[thick] (Aib)--(Aob);
  \draw[thick] (Aib)--(Bib);
  \draw[thick] (Aib)--(Ait);
  \draw[thick] (Aob)--(Aot);
  \draw[thick] (Bib)--(Bit);

  % Centro de curvatura y radios auxiliares: hacen visible la geometría cilíndrica.
  \fill (Ov) circle (1pt);
  \draw[auxiliar] (Ov)--(Aob);
  \draw[auxiliar] (Ov)--(Bob);

  % Las tres longitudes diferenciales del elemento cilíndrico.
  \draw[cota] (Aib)--(Aob)
    node[midway,below right=-1pt] {$dr'$};
  \coordinate (ArcStart) at
    ({\cx+\ex+1.52*cos(\anga)},{\cy+\ey+1.52*sin(\anga)});
  \draw[cota] (ArcStart)
    arc[start angle=\anga,end angle=\angb,radius=1.52]
    node[midway,right=2pt] {$r'\,d\varphi'$};
  \draw[cota] (Bob)--(Bot)
    node[midway,right=2pt] {$dz'$};

  \node[font=\bfseries] at (0.65,1.82)
    {Elemento cilíndrico $dV'$};
  \node at (0.65,-1.18)
    {$dV'=r'\,dr'\,d\varphi'\,dz'$};
\end{scope}

% Llamada desde el punto fuente a la cuña anular ampliada.
\draw[auxiliar,-{Stealth[length=1.5mm]}] (Q)--(6.45,3.65);
\end{tikzpicture}
\caption{Esquema tridimensional para la integración directa. El tramado marca
la región cargada $a<r'<b$; $\vec{x}$ y $\vec{x}'$ denotan, respectivamente,
los puntos de observación y fuente. A la derecha se amplía como cuña anular el
elemento de volumen cilíndrico asociado al punto $Q$.}
\label{fig:integracion-directa}
\end{figure}

Por simetría, la componente $\hat{y}$ se anula al integrar en $\varphi'$ y la
componente $\hat{z}$ se anula por antisimetría al integrar en $z'$. Solo
sobrevive la componente en dirección $\hat{r}$, que coincide con $\hat{x}$ en
el punto de evaluación:
\begin{equation}
E(r)=\frac{1}{4\pi\varepsilon_0}
\int_a^b\frac{\rho_0}{r'}r'\,dr'
\int_0^{2\pi}d\varphi'
\int_{-\infty}^{\infty}dz'\,
\frac{r-r'\cos\varphi'}
{\left[(r-r'\cos\varphi')^2+(r'\sin\varphi')^2+z'^2\right]^{3/2}}.
\end{equation}

Para efectuar primero la integral en $z'$, definimos
\[
R^2=(r-r'\cos\varphi')^2+(r'\sin\varphi')^2
=r^2+r'^2-2rr'\cos\varphi'.
\]
La antiderivada requerida es
\[
\int\frac{dz'}{(R^2+z'^2)^{3/2}}
=\frac{z'}{R^2\sqrt{R^2+z'^2}}+C.
\]
Al evaluar en los límites impropios se obtiene
\begin{equation}
\begin{aligned}
\int_{-\infty}^{\infty}\frac{dz'}{(R^2+z'^2)^{3/2}}
&=\left[
\frac{z'}{R^2\sqrt{R^2+z'^2}}
\right]_{-\infty}^{\infty}\\
&=\frac{1}{R^2}-\left(-\frac{1}{R^2}\right)
=\frac{2}{R^2}.
\end{aligned}
\end{equation}
Así, el campo queda reducido a la integral doble
\begin{equation}
E(r)=\frac{\rho_0}{2\pi\varepsilon_0}
\int_a^b dr'\int_0^{2\pi}d\varphi'\,
\frac{r-r'\cos\varphi'}
{r^2+r'^2-2rr'\cos\varphi'}.
\label{eq:integral-doble}
\end{equation}

Mediante el uso de identidades de potencial logarítmico o integrales de línea asociadas a potenciales de línea infinitos, el resultado para el campo eléctrico en las tres regiones es:
\[
E(r)=
\begin{cases}
0, & 0\leq r<a,\\[4pt]
\dfrac{\rho_0}{\varepsilon_0 r}\displaystyle\int_a^r dr'
=\dfrac{\rho_0}{\varepsilon_0}\left(1-\dfrac{a}{r}\right),
& a<r<b,\\[10pt]
\dfrac{\rho_0}{\varepsilon_0 r}\displaystyle\int_a^b dr'
=\dfrac{\rho_0}{\varepsilon_0}\left(\dfrac{b-a}{r}\right),
& r>b.
\end{cases}
\]
Finalmente, restituyendo la dirección radial, obtenemos
\begin{equation}
\boxed{
\vec E(r)=
\begin{cases}
\vec 0, & 0\leq r<a,\\[4pt]
\dfrac{\rho_0}{\varepsilon_0}
\left(1-\dfrac{a}{r}\right)\hat{r}, & a<r<b,\\[9pt]
\dfrac{\rho_0}{\varepsilon_0}
\left(\dfrac{b-a}{r}\right)\hat{r}, & r>b.
\end{cases}}
\label{eq:campo-regiones}
\end{equation}

\inciso{b} Para verificar el resultado mediante la ley de Gauss elegimos una
superficie gaussiana
\emph{cilíndrica}, coaxial con la distribución, de radio $r$ y longitud $h$. No se
trata de una esfera: una superficie esférica no respetaría la simetría de traslación
del problema.

La figura~\ref{fig:gauss} muestra cortes transversales de las superficies cilíndricas
gaussianas para las tres regiones. La trama oblicua identifica la distribución de
carga; al superponer una segunda familia de líneas se distingue la porción encerrada
por la superficie gaussiana.

\begin{figure}[htbp]
\centering
\begin{tikzpicture}[
    scale=0.90,
    transform shape,
    line cap=round,line join=round,
    gaussiana/.style={black,very thick,densely dashed},
    radio/.style={{Stealth[length=1.6mm]}-{Stealth[length=1.6mm]}},
    every node/.style={font=\small}
]
\def\ra{0.90}
\def\rb{1.62}

% Macro: centro, radio gaussiano, título, carga encerrada, etiqueta y trama extra.
\newcommand{\panelgauss}[6]{%
  \begin{scope}[shift={(#1,0)}]
    % Distribución completa: una sola familia de líneas.
    \path[
      pattern={Lines[angle=45,distance=4.0pt,line width=0.22pt]},
      pattern color=black!48,
      even odd rule
    ]
      (0,0) circle (\rb) (0,0) circle (\ra);

    % La trama adicional, inclinada en sentido opuesto, marca la carga encerrada.
    #6

    % Fronteras físicas, superficie gaussiana y eje común.
    \draw[thick] (0,0) circle (\rb);
    \draw[thick] (0,0) circle (\ra);
    \draw[gaussiana] (0,0) circle (#2);
    \fill (0,0) circle (1.3pt);
    \node[above left=-1pt] at (0,0) {$z\,\odot$};
    \draw[-{Stealth[length=1.7mm]}]
      (0,0)--({#2*cos(18)},{#2*sin(18)})
      node[pos=.58,above] {$r$};
    \draw[-{Stealth[length=1.5mm]}]
      (0,0)--({\ra*cos(218)},{\ra*sin(218)})
      node[pos=.66,above left=-1pt] {$a$};
    \draw[-{Stealth[length=1.5mm]}]
      (0,0)--({\rb*cos(292)},{\rb*sin(292)})
      node[pos=.72,below right=-1pt] {$b$};
    \node[font=\bfseries] at (0,2.34) {#3};
    \node[align=center,font=\scriptsize] at (0,-2.18) {#4};
    \node[font=\scriptsize,fill=white,inner sep=1pt]
      at #5 {$S_r=\partial V_r$};
  \end{scope}
}

% Región 1: no hay carga encerrada.
\panelgauss{0}{0.58}{Región 1: $r<a$}
  {$q_{V_r}=0$}
  {(1.12,0.66)}{}

% Región 2: trama cruzada en a<r'<r.
\panelgauss{5.15}{1.28}{Región 2: $a<r<b$}
  {$q_{V_r}=2\pi h\rho_0(r-a)$}
  {(1.34,1.10)}
  {\path[
      pattern={Lines[angle=135,distance=3.2pt,line width=0.25pt]},
      pattern color=black!72,
      even odd rule
    ]
     (0,0) circle (1.28) (0,0) circle (\ra);}

% Región 3: toda la distribución queda encerrada.
\panelgauss{10.30}{2.02}{Región 3: $r>b$}
  {$q_{V_r}=2\pi h\rho_0(b-a)$}
  {(1.66,1.48)}
  {\path[
      pattern={Lines[angle=135,distance=3.2pt,line width=0.25pt]},
      pattern color=black!72,
      even odd rule
    ]
     (0,0) circle (\rb) (0,0) circle (\ra);}

% Leyenda común.
\begin{scope}[shift={(5.15,-3.05)}]
  \draw[gaussiana] (-4.35,0)--(-3.55,0);
  \node[anchor=west,font=\scriptsize] at (-3.40,0)
    {superficie cilíndrica gaussiana de longitud $h$};
  \path[
    pattern={Lines[angle=45,distance=3.0pt,line width=0.22pt]},
    pattern color=black!48
  ] (-4.35,-0.48) rectangle (-3.95,-0.24);
  \draw[thin] (-4.35,-0.48) rectangle (-3.95,-0.24);
  \node[anchor=west,font=\scriptsize] at (-3.80,-0.36)
    {distribución de carga};
  \path[
    pattern={Lines[angle=45,distance=3.0pt,line width=0.22pt]},
    pattern color=black!48
  ] (0.02,-0.48) rectangle (0.42,-0.24);
  \path[
    pattern={Lines[angle=135,distance=2.7pt,line width=0.24pt]},
    pattern color=black!72
  ] (0.02,-0.48) rectangle (0.42,-0.24);
  \draw[thin] (0.02,-0.48) rectangle (0.42,-0.24);
  \node[anchor=west,font=\scriptsize] at (0.57,-0.36)
    {porción encerrada};
\end{scope}
\end{tikzpicture}
\caption{Ley de Gauss en las tres regiones. Cada panel es un corte transversal
de cilindros coaxiales. La línea discontinua representa $S_r=\partial V_r$, de
radio $r$ y longitud $h$; la trama cruzada marca la carga $q_{V_r}$ encerrada.}
\label{fig:gauss}
\end{figure}

La ley de Gauss en forma integral establece
\begin{equation}
\oint_{S_r}\vec E\cdot d\vec S
=\frac{q_{V_r}}{\varepsilon_0},
\qquad S_r=\partial V_r.
\end{equation}
El flujo a través de las tapas superior e inferior es nulo, pues allí
$\vec E\perp d\vec S$. Sobre la superficie lateral, en cambio, $\vec E$ es
paralela a $d\vec S$ y su magnitud es constante. Por tanto,
\begin{equation}
E(r)(2\pi r h)=\frac{q_{V_r}}{\varepsilon_0}.
\end{equation}

Para $r<a$ no hay carga encerrada, de modo que
\begin{equation}
q_{V_r}=0
\quad\Longrightarrow\quad E(r)=0.
\end{equation}

Para $a<r<b$, la carga encerrada abarca desde el radio interior $a$ hasta el
radio de la superficie gaussiana:
\begin{equation}
\begin{aligned}
q_{V_r}
&=\int_0^h dz'\int_a^r dr'\int_0^{2\pi}d\varphi'\,
  r'\left(\frac{\rho_0}{r'}\right)\\
&=2\pi h\rho_0\int_a^r dr'
 =2\pi h\rho_0(r-a).
\end{aligned}
\end{equation}
Sustituyendo en la ley de Gauss, tenemos
\begin{equation}
E(r)=\frac{2\pi h\rho_0(r-a)}
{\varepsilon_0(2\pi r h)}
=\frac{\rho_0}{\varepsilon_0}\left(1-\frac{a}{r}\right).
\end{equation}

Para $r>b$ queda encerrada toda la carga del cascarón:
\begin{equation}
q_{V_r}=2\pi h\rho_0(b-a).
\end{equation}
Por tanto,
\begin{equation}
E(r)=\frac{2\pi h\rho_0(b-a)}
{\varepsilon_0(2\pi r h)}
=\frac{\rho_0}{\varepsilon_0}\left(\frac{b-a}{r}\right).
\end{equation}

El resultado coincide con la ecuación~\eqref{eq:campo-regiones} obtenida mediante
integración directa.

\inciso{c} Para calcular el potencial usamos
\begin{equation}
\phi(r)-\phi(r_{\mathrm{ref}})
=-\int_{r_{\mathrm{ref}}}^{r}E(u)\,du,
\end{equation}
donde se ha usado $u$ como variable muda para no confundirla con el radio del punto
de observación. Esta es la reducción radial de
$\phi(\vec{x})=\phi(\vec{x}_0)-\int_{\vec{x}_0}^{\vec{x}}\vec E\cdot d\vec{x}$.
Por enunciado, la condición de referencia es $\phi(b)=0$.

Para $r>b$ integramos desde el punto de referencia $r=b$:
\begin{equation}
\phi(r)=-\int_b^r\frac{\rho_0(b-a)}{\varepsilon_0u}\,du
=-\frac{\rho_0(b-a)}{\varepsilon_0}\ln\left(\frac{r}{b}\right).
\end{equation}

Para $a<r<b$ partimos nuevamente desde $r=b$, ahora hacia el interior:
\begin{equation}
\begin{aligned}
\phi(r)
&=-\int_b^r\frac{\rho_0}{\varepsilon_0}
\left(1-\frac{a}{u}\right)du\\
&=-\frac{\rho_0}{\varepsilon_0}
\left[(r-b)-a\ln\left(\frac{r}{b}\right)\right]\\
&=\frac{\rho_0}{\varepsilon_0}
\left[(b-r)+a\ln\left(\frac{r}{b}\right)\right].
\end{aligned}
\end{equation}

Para $r<a$ el campo es cero. Por esto el potencial es constante e igual al
valor en la frontera interior $r=a$:
\begin{equation}
\phi(r)=\phi(a)=\frac{\rho_0}{\varepsilon_0}
\left[(b-a)+a\ln\left(\frac{a}{b}\right)\right].
\end{equation}

En conjunto,
\begin{equation}
\boxed{
\phi(r)=\frac{\rho_0}{\varepsilon_0}
\begin{cases}
(b-a)+a\ln\left(\dfrac{a}{b}\right), & 0\leq r<a,\\[7pt]
(b-r)+a\ln\left(\dfrac{r}{b}\right), & a<r<b,\\[7pt]
-(b-a)\ln\left(\dfrac{r}{b}\right), & r>b.
\end{cases}}
\end{equation}

\clearpage
\section*{Visualización numérica}

Para comparar las tres regiones tomamos $a=L$ y $b=2L$, con $L$ una longitud
de referencia, y definimos
\[
\widetilde{\rho}=\frac{\rho L}{\rho_0},
\qquad
\widetilde E=\frac{\varepsilon_0E}{\rho_0},
\qquad
\widetilde\phi=\frac{\varepsilon_0\phi}{\rho_0L}.
\]
La figura~\ref{fig:visualizacion-cilindro} muestra los perfiles radiales que se
obtienen de las expresiones analíticas anteriores. El corte transversal permite
ver que el campo se anula en la cavidad, crece dentro de la región cargada y
decrece como $1/r$ en el exterior.

\begin{figure}[htbp]
\centering
\includevisualizacion[width=0.97\textwidth]{cilindro_hueco_visualizacion}
\caption{Visualización numérica para $a=L$, $b=2L$ y $\rho_0>0$. Los paneles
(a)--(c) muestran la densidad, el campo y el potencial adimensionales. En (d),
el color identifica la densidad volumétrica en el cascarón y las flechas
representan el campo eléctrico sobre un corte perpendicular al eje.}
\label{fig:visualizacion-cilindro}
\end{figure}

\end{document}
